\chapter*{Abstract}
\addcontentsline{toc}{chapter}{Abstract}

Large-scale connectomics --- the mapping of neural circuits from volumetric
electron microscopy (EM) and X-ray nano-holotomography (XNH) data --- requires
automated segmentation of three-dimensional image volumes into individual neuronal
fragments. State-of-the-art automated segmenters produce fragments with high accuracy
but inevitably introduce \emph{split errors}, where a single neuron is represented as
multiple disconnected fragments. Correcting these errors at the scale of modern
connectomics datasets (tens to hundreds of teravoxels) demands automated
proofreading pipelines that can detect and propose corrections for split errors
with high recall and with explicit handling of uncertainty.

This report presents a modular, conservative post-segmentation pipeline that
takes the output of an automated segmenter and produces a graph-based
connectivity representation with confidence-annotated merge proposals. The pipeline
operates in six stages: fragment extraction with TEASAR skeletonization, graph
construction via a skeleton-node KD-tree index, candidate generation with a
four-factor composite scoring model, conservative three-outcome rule-based
validation (ACCEPT / REJECT / AMBIGUOUS), structure assembly, and multi-format
export. Critically, the pipeline treats uncertainty as a first-class output ---
ambiguous connections are preserved for human review rather than forced to a
binary decision.

Two key innovations drive the approach's effectiveness on axon-dominated datasets.
First, a \emph{skeleton-node graph} indexes every skeleton node from every fragment
(rather than only degree-1 endpoints), exposing interior splits along long axons
that an endpoint-only graph cannot detect. Second, \emph{distance-conditioned
scoring} switches from a proximity-weighted composite to a proximity-free,
alignment-continuity-weighted score for long-range candidate pairs, preventing
the artificial suppression of genuine long-range splits.

Evaluated on the CREMI Drosophila EM benchmark, the pipeline achieves
Precision = 1.000, Recall = 0.909, and F1 = 0.952. On the XPRESS mouse white
matter XNH challenge --- the primary domain-appropriate benchmark --- the pipeline
achieves Coverage = 83.5\% and Recall = 0.994 on the training volume (Experiment~17),
and Coverage = 72.4\%, Recall = 0.885 on the fully held-out validation volume
(Experiment~23), confirming that the learned thresholds generalize without
overfitting. The pipeline is validated by 428 automated tests and includes a
reproducible visual diagnostic report for per-candidate inspection.

The main open challenge is precision: approximately 354{,}000 false-positive
merge proposals are generated alongside 1{,}245 true positives at XPRESS scale
(Precision $\approx$ 0.004), reflecting the difficulty of discriminating
same-axon from different-axon long-range pairs in myelinated white matter using
geometric features alone. Future work targets this gap through improved scoring
and lightweight machine-learned filters.
